\chapter{Introduction and motivation}\label{ch:intro}

\section{Generation, ancestry, and what an observation can reveal}
Consider a rule that builds objects step by step. Starting from a seed, each step applies one of finitely many operations, an operation may be available only when the current object satisfies a condition, and the order of operations may or may not matter. A \emph{history} is the record of seed and operations; the object it produces is its \emph{endpoint}. In almost every interesting system a single endpoint $t$ arises from many histories. The set of those histories is the \emph{fibre} $F_t$ of $t$, and the base-two logarithm of its size, $A(t)=\log_2|F_t|$, measures the \emph{ancestral ambiguity} of $t$: the number of bits that remain unknown about how $t$ was made when only $t$ is seen.

A single number does not describe the problem. Some histories are related by harmless reorderings of independent steps; others differ in seed, in the multiset of operations used, or in the \emph{constraints} that the grammar imposes on order. Two systems may have exactly as many histories per endpoint and still differ in whether their histories form a connected, hole-free family. This monograph develops the mathematics needed to say precisely what is and is not recoverable about generative ancestry from observations of increasing strength: the endpoint, the number of histories, local statistics of the reordering structure, and the global topology of the family of histories.

The setting in which this can be done exactly is the following. Fix a finite set of $k$ distinct events (labels) and a set $A$ of orderings of them, the \emph{admissible histories}. The orderings are the vertices of the permutohedron $P_k$; two orderings are adjacent when they differ by exchanging two neighbouring events (a \emph{legal swap}). The faces of $P_k$ whose vertices all lie in $A$ form a complex $K_A$, and the homology of $K_A$ records whether the admissible histories hang together without holes: a one-dimensional class is a loop of histories that cannot be contracted by commuting independent events, a two-dimensional class is a hollow shell, and so on. For arithmetic rules the answer depends on the divisibility of the seed, so the topology can switch on and off along a divisibility order.

\section{Five questions}
The monograph is organized around five questions, each of which has an exact answer in the framework, and each of which is answered below by theorems, counterexamples or explicitly stated open problems.
\begin{enumerate}[label=\textbf{Q\arabic*.},leftmargin=2.6em]
\item \emph{Multiplicity.} How many histories reach an endpoint, how does that number grow, and how does the ambiguity decompose into source, content, reordering class and order within a class? (Chapter~\ref{ch:ancestry}.)
\item \emph{Admissibility.} Which sets of orderings can a given rule produce, how do they depend on the seed, and what combinatorial structure does the dependence have? (Chapters~\ref{ch:arith}--\ref{ch:swaps}.)
\item \emph{Topology.} When do admissible sets have holes in dimension two and higher, when do holes survive as the seed is made more divisible, and how are holes of dimension $d{+}1$ built from holes of dimension $d$? (Chapters~\ref{ch:topology}--\ref{ch:dimension}.)
\item \emph{Recoverability.} Which observations of a system---history count, local reordering statistics, ball sizes in the swap graph---determine whether a hole is present, and which do not? (Chapters~\ref{ch:access}--\ref{ch:predictive}.)
\item \emph{Forbidden histories.} Can topology be read from what is forbidden rather than what is allowed, how small a part of the forbidden structure decides it, and does the answer depend on integer arithmetic or only on the form of the constraints? (Chapters~\ref{ch:duality}--\ref{ch:locality}.)
\end{enumerate}

\section{Summary of results}\label{sec:summary}
We list the principal results with pointers. Each is stated with hypotheses and a status in the cited place; here the status tags are \textsf{A} (analytic), \textsf{CA} (computer-assisted), \textsf{FE} (finite-exact), \textsf{B} (bounded), \textsf{F} (falsified).
\begin{enumerate}[label=\textbf{R\arabic*.},leftmargin=2.6em]
\item \textbf{Ancestry and its decomposition} (\textsf{A}). Fibres are finite, complete histories exist for every endpoint, and the ambiguity decomposes exactly as $A(t)=\Ent(M|T)+\Ent(Q|M,T)+\Ent(C|M,Q,T)+\Ent(W|C,T)$ under the uniform ensemble on the fibre (\cref{thm:hierarchy}). The endpoint $72$ has four histories and decomposes as $2=1.5+0.5$ bits (\cref{ex:72}); $|F_{72^r}|\ge\binom{2r}{r}$, with exact counts through $r=10$ (\cref{prop:72r}).
\item \textbf{Prefix-defined admissibility} (\textsf{A}). For the arithmetic rule, a history is admissible iff the denominator of every prefix product divides the seed (\cref{lem:F}); the admissible sets form a lcm-closure system with unique minimal seeds, and nested complexes along divisibility (\cref{thm:closure,thm:filtration}). Block factorization holds for \emph{every} prefix-defined system and rules out six-vertex slices of commuting cubes (\cref{thm:block}).
\item \textbf{Reconstruction from the swap graph} (\textsf{A}, \textsf{CA}). The complex $K_A$ is determined by the generator-coloured swap graph for arbitrary $A$ (\cref{thm:swap-reconstruction}); for prefix-defined $A$ the uncoloured graph determines the two-skeleton (\cref{thm:skeleton}), and face recognition is sound for convex $A$ in the verified range (\cref{thm:face-recognition}). Without the cube-product property the graph alone fails (\cref{prop:petrie}).
\item \textbf{Survival of two-dimensional holes} (\textsf{A}, \textsf{FE}). An explicit 96-face two-cycle and six-face cocycle survive precisely when none of the four minimal filling families of prefix sets is admissible (\cref{thm:anchor,thm:hole-survival,thm:exact-survival}); the necessity of two of the three single-gate conditions, conjectured at first, is refuted by explicit examples (\cref{prop:gate-nec}).
\item \textbf{Lifting holes to higher dimension} (\textsf{A}, \textsf{FE}, \textsf{B}). A suspension-type operator raises cycles from $P_k$ to $P_{k+1}$; its boundary identity, the cycle property, and the vanishing of diagonal lifts are proved; certified three- and four-dimensional arithmetic examples exist; and the lifted class depends on the chosen filling through an exactly characterised map (\cref{lem:LT0,thm:lift-cycle,thm:filling-dependence,thm:iterated-lift,thm:JX}).
\item \textbf{Entropy does not see topology} (\textsf{FE}, \textsf{A}). Systems with equal history counts, equal entropies and equal radius-one statistics can differ in $\beta_2$; the exact radius-two identity pins down what radius-two data contain (\cref{thm:radius2,prop:witnesses,obs:equal-N}). Pre-registered predictions were partly falsified; plug-in information overstates predictive information (\cref{obs:preregistered,obs:S2-falsified,obs:oos}).
\item \textbf{Forbidden-history duality and the dead-set nerve} (\textsf{A}). $\Hh_j(K_A)\cong\Hh^{k-3-j}(Y_A)$; for prefix-defined $A$, $Y_A\simeq N[D_A]$ and $\beta_2(K_A)=\dim H_1(N[D_A])$ for $k=6$ (\cref{thm:AD,lem:star,thm:nerve}).
\item \textbf{Minimal structural information} (\textsf{FE}). Over a six-label census of $219{,}145$ dead-set patterns, no tested count-level statistic of the nerve determines $\beta_2>0$; the strong-collapse core does (\cref{obs:collapse}).
\item \textbf{Locality of constraints} (\textsf{CA}, \textsf{FE}). Pairwise-adjacency grammars on six events have $\beta_2=0$ (\cref{thm:pairwise}); two forbidden consecutive triples give $\beta=(1,2,5,0,0,0)$ (\cref{prop:triple}). No claim is made for arbitrary grammar length.
\end{enumerate}

\section{The central thesis and what is not claimed}
The thesis of the monograph is that \emph{generative information has more than one layer}: multiplicity and probability of histories (Part~I), arithmetic or symbolic constraints that determine admissibility (Part~II), topology of the family of histories (Parts~III), forbidden-history structure (Part~V), and the observations that can or cannot recover each of these (Part~IV). Part~VI relates the layers and states what remains open.

We do not claim that the SELECT-II arithmetic rule is the only legitimate generative rule of the UNNS substrate; it is the first rule in which all layers can be made exact and it serves as the principal worked realization. We do not claim a replacement for Shannon's theory, a universal physical entropy, or any result for generative evolution of collective fields (that remains \cref{prob:evolution}). We do not claim that statements verified on a census hold beyond it. Where a result is a corpus-scoped observation, it is labelled so and is never promoted to a theorem.

\section{Conventions, status vocabulary and verification}\label{sec:conventions}
\paragraph{Coefficients and arithmetic.} Homology and cohomology have coefficients in $\F=\mathrm{GF}(2)$ unless stated. All arithmetic in computations is exact: integers, rational numbers (\texttt{Fraction}) and linear algebra over $\mathrm{GF}(2)$ on bit-vectors; no floating-point value enters any equality. Entropies are reported to the stated number of digits and are derived from exact integer counts.
\paragraph{Status vocabulary.}
\begin{description}[leftmargin=1.2em,labelsep=0.5em,itemsep=0.1em]
\item[Analytic] a complete human-checkable proof is given in the text (classical theorems used are named and cited).
\item[Computer-assisted] the proof reduces to a stated finite family of exact statements; the finite statements are certified by named certificates or checkers and, where the manuscript's own verifier re-derives them, this is stated. The checker is short enough to be read.
\item[Finite-exact] an exact computation over a stated finite domain. No claim beyond the domain is made, and the domain is always stated.
\item[Bounded] holds on a searched or sampled domain; exhaustive only within the stated bounds; not a theorem.
\item[Conjectural / open] stated as a conjecture or problem; no proof is claimed.
\item[Falsified] a stated conjecture or pre-registered prediction that was refuted; we keep it and the counterexample.
\end{description}
\paragraph{Re-run policy.} Expensive inherited computations (the six-label census of $346{,}859$ source regimes, the seven-label scan of $1{,}560{,}780$ contents, the search for the first four-dimensional example, and similar) are recorded as \emph{not re-run}. For them we cite the certificate or register entry that carries the evidence. What the manuscript's own verifier \emph{does} re-derive independently---from the definitions and without importing the earlier engines---is listed in Appendix~\ref{app:methods}.

\section{Provenance of the material}\label{sec:provenance}
The monograph integrates the Generative Information Geometry programme of the UNNS (Unbounded Nested Number Sequences) Research Collective: the seven releases GIG-I v001--v007 (the fifth in three stages) and the six releases GIG-INFO v001--v006, together with the focused paper on forbidden histories. The arithmetic realization is the frozen UNNS operator SELECT-II, $T_N(m)=m(N{+}1)^2/N$. The source releases are frozen: this monograph reads them and does not modify them. Each theorem in the text names its release in a \emph{Provenance} line; \cref{app:sources} indexes the releases and the theorem crosswalk.

Two provenance points are disclosed. First, an error in the second-level build of one release (the fifth GIG-INFO release recorded the same parent hash for two different parent releases because of a path substitution) was found and documented in an erratum; no result depended on it (\cref{app:sources}). Second, no pinned commit hash of the project repository was available when this monograph was prepared, so the bibliography and the source index list release folders and manifest hashes but do not print repository links; adding links is a matter of recording a commit once one is fixed.

\section{Reading guide}
\Cref{fig:dependencies} shows the logical dependence of the parts. Readers interested only in the duality and the nerve theorem may read Chapters~\ref{ch:ancestry} (up to the definition of $K_A$), \ref{ch:swaps}, \ref{ch:duality}--\ref{ch:locality}. Readers interested in arithmetic should read Part~II before Part~III. Chapter~\ref{ch:synthesis} contrasts the layers and Chapter~\ref{ch:open} lists open problems; neither is needed for the proofs.
\begin{figure}[ht]\centering
\begin{tikzpicture}[node distance=4mm,>=Stealth,every node/.style={font=\small}]
\node[draw,rounded corners,minimum width=2.7cm,align=center] (I) {I\\Foundations};
\node[draw,rounded corners,minimum width=2.7cm,align=center,right=9mm of I] (II) {II\\Arithmetic\\admissibility};
\node[draw,rounded corners,minimum width=2.7cm,align=center,right=9mm of II] (III) {III\\Higher\\topology};
\node[draw,rounded corners,minimum width=2.7cm,align=center,below=7mm of II] (IV) {IV\\Information\\recovery};
\node[draw,rounded corners,minimum width=2.7cm,align=center,right=9mm of IV] (V) {V\\Forbidden\\histories};
\node[draw,rounded corners,minimum width=2.7cm,align=center,below=7mm of IV] (VI) {VI\\Synthesis};
\draw[->] (I)--(II);\draw[->] (II)--(III);\draw[->] (I)--(IV);\draw[->] (II)--(IV);\draw[->] (III)--(V);\draw[->] (II)--(V);\draw[->] (IV)--(VI);\draw[->] (V)--(VI);\draw[->] (V)--(IV);
\end{tikzpicture}
\caption{Dependence of the six parts. The arrow from Part~V to Part~IV records that the explanation of what the observations of Part~IV can recover uses the duality of Part~V.}\label{fig:dependencies}
\end{figure}
